\documentclass{article}
\usepackage{graphicx} % Required for inserting images

\usepackage{graphicx} % Required for inserting images
\usepackage{graphicx} % Required for inserting images
\usepackage[utf8]{inputenc} %for making sure that...
\usepackage[T1]{fontenc} %Swedish characters (å,ä,ö) appear
\usepackage{amsmath,amsthm,amsfonts,amssymb,mathrsfs} %this set of packages allows us to use basically any mathematical notation
\usepackage{comment} %to comment out part of latex using \begin{comment}
\usepackage{physics} %for \abs{}
%graphics related
\usepackage{enumitem}
\usepackage{graphicx} %in particular, this package allows us to insert nice figures and captions
\usepackage{subcaption}
\usepackage[export]{adjustbox}
\usepackage{wrapfig}
%The following is to set up listings with coding styles
\usepackage{xcolor,color,listings}
%New colors defined below
\definecolor{codegreen}{rgb}{0,0.6,0}
\definecolor{codegray}{rgb}{0.5,0.5,0.5}
\definecolor{codepurple}{rgb}{0.58,0,0.82}
\definecolor{backcolour}{rgb}{0.95,0.95,0.92}
%Code listing style named "mystyle"
\lstdefinestyle{mystyle}{
  backgroundcolor=\color{backcolour},   commentstyle=\color{codegreen},
  keywordstyle=\color{magenta},
  numberstyle=\tiny\color{codegray},
  stringstyle=\color{codepurple},
  basicstyle=\ttfamily\footnotesize,
  breakatwhitespace=false,         
  breaklines=true,                 
  captionpos=b,                    
  keepspaces=true,                 
  numbers=left,                    
  numbersep=5pt,                  
  showspaces=false,                
  showstringspaces=false,
  showtabs=false,                  
  tabsize=2
}
%"mystyle" code listing set
\lstset{style=mystyle}

\usepackage[a4paper, total={6in, 8in}]{geometry}
%so that we can pleasantly insert webadresses
\usepackage{hyperref}
%to make it shorter to write these sets of numbers
\newcommand{\R}{\mathbb{R}}
\newcommand{\Z}{\mathbb{Z}}
\newcommand{\N}{\mathbb{N}}
\newcommand{\Q}{\mathbb{Q}}
\newcommand{\C}{\mathbb{C}}

%To get white-space parapgraphs
\parskip6pt
\parindent0pt
\usepackage{tikz}
\usepackage{amsmath}
\newcommand{\hly}[1]{\colorbox{lightgray}{\parbox{\textwidth}{#1}}}
\usetikzlibrary{arrows}

\author{\vspace{-5ex}}
\date{\vspace{-5ex}}
\title{NUMB11 Assignment2}
\author{Sek Huen Leung, Nancy Truong}
\date{September 2024}

\begin{document}

\maketitle

\section{Task 1}
In this task, we determine the coefficients from a measurement modelling a signal at 200 equidistant sampling points in $[0, 4\pi]$. QR-factorisation and SVD method were utilized and compared with respect to their computational time. QR-factorisation has a shorter computational time than SVD. And the signal was also plotted. See Figure 1-3 below.

\begin{figure}
    \centering
    \includegraphics[width=1\linewidth]{NUM11_assignment3_task1.png}
    \caption{Task 1}
    \label{fig:enter-label}
\end{figure}

\begin{figure}
    \centering
    \includegraphics[width=1\linewidth]{NUMB11_assignment3_task1b.png}
    \caption{Task 1}
    \label{fig:enter-label}
\end{figure}

\begin{figure}
    \centering
    \includegraphics[width=1\linewidth]{NUMB11_assignment3_task1c.png}
    \caption{Task 1}
    \label{fig:enter-label}
\end{figure}
\newpage 
\section{Task 2}
In this task, we choose $b$ and $\delta b$ such that 
$$\frac{\norm{\delta x}_2}{\norm{x}_2}= \kappa_2 (A) \frac{\norm{\delta b}_2}{\norm{b}_2}$$ when solving $Ax=b$.\\
We know that 
$$\kappa_2(A) = \norm{A}_2 \norm{A^{-1}}_2 = \frac{\sigma_1}{\sigma_n}$$
where $\sigma_1$ is the largest singular value of $A$ and $\sigma_n$ is the smallest singular value.\\
To satisfy the equality, we choose
$$b = u_1 \text{ and } \delta b = u_n,$$
where $u_1$ and $u_n$ are the left singular vectors corresponding to $\sigma_1$ and $\sigma_n$ respectively.\\
We solve
$$Ax = U \Sigma V^T x = b = u_1$$
we have
\begin{equation*}
    \begin{aligned}
        \Sigma V^T x &= U^T u_1 \\
        \Sigma V^T x &= e_1 \\
         V^T x &= \Sigma^{-1} e_1 \\
         x &=V \frac{e_1}{\sigma_1} \\
         x &= \frac{v_1}{\sigma_1}. \\
    \end{aligned}
\end{equation*}
Then we solve the perturbed equation
$$A(x+\delta x) = U \Sigma V^T (x+\delta x) = b + \delta b= u_1 + u_n$$
we have
\begin{equation*}
    \begin{aligned}
        \Sigma V^T (x + \delta x)  &= U^T (u_1 +u_n) \\
         x + \delta x &= \frac{v_1}{\sigma_1} + \frac{v_n}{\sigma_n} \\
         \delta x = \frac{v_n}{\sigma_n}.
    \end{aligned}
\end{equation*}
The left hand side of the equality
$$\frac{\norm{\delta x}_2}{\norm{x}_2} = \frac{\norm{\frac{v_n}{\sigma_n}}}{\norm{\frac{v_1}{\sigma_1}}} = \frac{\sigma_1}{\sigma_n},$$
and the right hand side 
$$\kappa_2 (A) \frac{\norm{\delta b}_2}{\norm{b}_2} = \frac{\sigma_1}{\sigma_n} \cdot \frac{1}{1} = \frac{\sigma_1}{\sigma_n}.$$
\qed
\newpage 
\section{Task 3}
In this task, we write a program to verify the equality in Task 2. We tried different method, the codes and results are shown in \textit{Figure 4} and \textit{Figure 5}. We got different results from different method and we think Hilbert matrix is so ill that we got error from calculating the error.
\begin{figure}
    \centering
    \includegraphics[width=0.9\linewidth]{A3T3A.png}
    \caption{Task 3 Code}
    \label{fig:enter-label}
\end{figure}
\begin{figure}
    \centering
    \includegraphics[width=1\linewidth]{A3T3B.png}
    \caption{Task 3 Result}
    \label{fig:enter-label}
\end{figure}
\newpage 
\section{Task 4}
By definition, a $n \cross n$ square matrix $[A]$ is a strictly diagonally dominant matrix if
$$
\abs{a_{ii}} > \sum_{j=1, i\neq j} \abs{a_ij}
$$
for $i=1,2,...,n$.
That is, for each row, the absolute value of the diagonal element is strictly greater than the sum of the absolute values of the rest of the elements of the that row.

We assume there exists a vector $u \neq 0$ such that $Au=0$. This means that for $i=1,2,...,n,$


$$ 
\sum_{j=1}^n a_{ij}u_j = 0
$$

Let $u_m = \norm{u}_\infty \neq 0$ meaning that $u_m$ is the largest entry of vector $u$ by absolute value.

We have
\begin{equation*}
    \begin{aligned}
    & &0 &= \sum_{j=1}^n a_{ij}u_j \\
    &\iff &a_{mm}u_m &= - \sum_{j\neq m}^n a_{ij}u_j \\
    &\iff &a_{m} &= - \sum_{j\neq m}^n a_{ij} \frac{m_j}{u_m} \\
    \end{aligned}
\end{equation*}

Taking the absolute value, we get
\begin{equation*}
    \begin{aligned}
        \abs{a_{mm}} &= \abs{\sum_{j\neq m}^n a_{ij} \frac{x_j}{x_m}}\\
        &\leq \sum_{j\neq m}^n \abs{a_{ij} \frac{x_j}{x_m}} \\
        &= \sum_{j\neq m}^n \abs{a_{mj}}\abs{\frac{x_j}{x_m}}
    \end{aligned}
\end{equation*}
Since $x_m > x_j$, $\frac{x_j}{x_m}<1$ for all $j\neq m$.

Therefore, 
$$
\abs{a_{mm}} < \sum_{j\neq m}^n \abs{a_{mj}}
$$
This contradicts the definition of a strictly diagonally dominant matrix.
\\
\\
Suppose $\abs{u_i} = \abs{u_j}$ for all $i,j = 1,2,...,n$. Let $\abs{u_i} = c$ for some $c>0$. 

\\
\\ 

We also have 
$$
\abs{a_{ii}u_i} =  \abs{\sum_{j\neq i}^n a_{ij}u_j} \leq \sum_{j\neq i}^n \abs{a_{ij}}\abs{u_j}
$$
Then
\begin{equation*}
    \begin{aligned}
        & &\abs{a_{ii}} \abs{u_i} &\leq \sum_{j\neq i}\abs{a_{ij}}\abs{u_j} \\
        &\iff &\abs{a_{ii}} c &\leq \sum_{j\neq i}\abs{a_{ij}}c \\
        &\iff &\abs{a_{ii}} &\leq \sum_{j\neq i}\abs{a_{ij}}
    \end{aligned}
\end{equation*}
This contradicts the assumption of values of all $\abs{u_j}$ being equal so that matrix $A$ is a strictly diagonally dominant matrix.


So $0 \not\in \ker(A)$. Therefore, A is invertible.


\section{Task 5}
We have
$$
u(x_i) \approx u_i
$$
$$
u'(x_i) \approx \frac{u_i-u_{i-1}}{\Delta x}
$$
\begin{equation*}
    \begin{aligned}
        u''(x_i) &\approx \left(\frac{u_{i+1}-u_1}{\Delta x} - \frac{u_i-u_{i-1}}{\Delta x}\right) \frac{1}{\Delta x} \\
        &= \frac{u_{i+1}-2u_i+u_{i-1}}{(\Delta x)^2} 
    \end{aligned}
\end{equation*}
So
\begin{equation*}
    \begin{aligned}
    & &-u''(x) &= f(x) \\
    &\iff &\frac{-u_{i-1}+2u_i-u_{i+1}}{(\Delta x)^2} &= f(x) \\
    &\iff &-u_{i-1}+2u_i-u_{i+1} &= f(x)(\Delta x)^2
    \end{aligned}
\end{equation*}
With the additional conditions $u(0) = u(1) = 0$, we can express this linear system in matrix form $Au = \textbf{f}$ such that
\begin{itemize}
    \item Vector u: $u=[u_1,u_2,...,u_{n-1}]$
    \item Matrix A of size $(n-1)\cross (n-1)$
        \begin{equation*}
            A =
            \begin{bmatrix}
                2 & -1 & 0 & \dots & 0 \\
                -1 & 2 & -1 & \dots & 0 \\
                0 & -1 & 2 & \dots & 0 \\
                \vdots & \vdots & \vdots & \ddots & \vdots \\
                0 & 0 & 0 & -1 & 2
            \end{bmatrix}
        \end{equation*}
    \item Vector \textbf{f}
    \begin{equation*}
    \textbf{f} = (\Delta x)^2 
        \begin{bmatrix}
        f(x_1) \\
        f(x_2) \\
        \vdots \\
        f(x_{n-1})
        \end{bmatrix}
    \end{equation*}
\end{itemize}

\\
\\

The matrix 
\begin{equation*}
    \begin{bmatrix}
        1 & -1 & 0 \\
        -1 & 2 & -1 & \ddots \\
        0 & \ddots & \ddots & \ddots & 0 \\
         & \ddots & -1 & 2 & -1 \\
        & & 0 & -1 & 1
    \end{bmatrix}
\end{equation*}
is nonsingular there is no a pair of columns or rows that are linearly dependent, thus, $\det(A) \neq 0$.
\end{document}
